% !TeX root = ../main.tex
\section{Ringkasan Verifikasi dan Batas Hasil}
\label{sec:verification}
Hasil lokal yang dilaporkan merupakan pengujian bertanggal \EvidenceDate. Verifikasi dibagi menjadi unit dan analisis statis dengan go vet, integrasi PostgreSQL, regresi stack, dan load test. Setiap jenis pengujian memberikan bukti dengan cakupan yang berbeda. Uji unit memeriksa logika tertentu, sehingga hasilnya tidak dapat langsung dipakai untuk menyimpulkan perilaku deployment, ketahanan mesin, atau kapasitas produksi.

\begin{tabularx}{\linewidth}{P{0.30\linewidth}Y}
\toprule Pemeriksaan & Hasil dan cakupan\\\midrule
Sembilan module Go & Bootstrap, unit test, dan vet; delapan service serta CLI.\\
PostgreSQL nyata & Transaksi ingest dan query lulus dalam schema test terisolasi.\\
Regresi stack & Fondasi, ingest, event, query, expiry token alami, serta rebuild independen lulus; durasi run 234,280 s.\\
Fresh subscriber & Group baru dan store kosong berhasil replay; suite event tambahan lulus dalam 71,575 s.\\
Readiness consumer & Test broker tidak merespons memenuhi deadline 50 ms dan menolak penumpukan probe.\\
Load final & Tiga skenario lulus; angka dan batas pengukuran pada Bagian~\ref{sec:p2}.\\
GitHub Actions & Workflow tersedia; hasil hosted belum dapat diverifikasi dari sesi penyusunan laporan.\\
\bottomrule
\end{tabularx}

Kegagalan awal tidak dihapus dari bukti. Regresi pertama menemukan timeout readiness consumer saat Kafka mati. Pengujian terisolasi menunjukkan \texttt{Ping} library dapat menunggu handshake sekitar 10 s meskipun context sudah deadline. Perbaikan membatasi waktu tunggu caller dan satu probe yang belum selesai; regresi penuh kemudian lulus. Hasil gagal pertama tersedia pada \evidence{regression-before-readiness-fix.txt}, sedangkan hasil akhir pada \evidence{regression.txt}.

\lstinputlisting[numbers=none,caption={Petikan keluaran pengujian aktual. Ini transkrip test, bukan screenshot yang disimulasikan.}]{assets/evidence/regression-excerpt.txt}

\pending{Status GitHub Actions dan tautan run final: \rule{0.53\linewidth}{0.3pt}. API tanpa autentikasi pada sesi penyusunan mengembalikan 404, sehingga keberhasilan CI hosted tidak diasumsikan.}

\subsection{Kesimpulan dan batas cakupan}
Alur ingest, query, otorisasi, serta distribusi event telah terhubung dalam implementasi. Bukti pengujian lokal mendukung skenario P1 hingga P5 yang dipetakan pada lampiran. Cakupan sistem masih terbatas pada PoC dalam satu host dan belum mencakup failover lintas node, koordinasi multi-instance Aggregator, limit terdistribusi, transport terenkripsi end-to-end, maupun pengiriman notifikasi nyata ke perangkat.

Demo sinkron tetap perlu dijalankan kelompok dan tidak digantikan oleh laporan test. Verifikasi panjang outage 10 hingga 20 menit secara langsung belum direpresentasikan oleh load test 20 detik yang disimpan; endpoint operator mendukung outage hingga dimatikan, tetapi kemampuan konfigurasi tidak sama dengan bukti durasi demo tertentu. Tag final dan pengumpulan juga merupakan tindakan terpisah setelah kelompok meninjau laporan.
